% SPDX-License-Identifier: CC-BY-4.0
\documentclass[11pt,letterpaper]{article}

\usepackage[T1]{fontenc}
\usepackage{amsmath,amssymb}
\usepackage{tikz}
\usepackage{pgfplots}
\pgfplotsset{compat=1.18}
\usepackage{algorithm}
\usepackage{algpseudocode}
\usepackage{booktabs}
\usepackage{enumitem}
\usepackage[hidelinks]{hyperref}

\newcommand{\N}{\mathbb{N}}
\newcommand{\Z}{\mathbb{Z}}
\newcommand{\Q}{\mathbb{Q}}
\newcommand{\R}{\mathbb{R}}
\newcommand{\C}{\mathbb{C}}
\newcommand{\Map}[3]{#1 \colon #2 \to #3}
\newcommand{\Mod}[1]{#1\text{-}\mathrm{Mod}}
\newcommand{\GroupHom}{\mathrm{Hom}_{\mathbf{Grp}}}
\newcommand{\RingHom}{\mathrm{Hom}_{\mathbf{Ring}}}

\newtheorem{theorem}{Theorem}
\newtheorem{lemma}[theorem]{Lemma}
\newtheorem{corollary}[theorem]{Corollary}
\newtheorem{definition}[theorem]{Definition}
\newtheorem{remark}[theorem]{Remark}
\newenvironment{proof}{\noindent\textit{Proof.}}{\hfill\(\square\)}

\title{Generalized Collatz Library}
\author{}
\date{\today}

\begin{document}
\maketitle
\section*{Signed Dyadic Affine Tails and Exact Cycles}

\begin{abstract}
For the division-first integer map with divisor two, opposite-parity
coefficients make the odd branch invariant. A translated integer coordinate
then gives a complete signed escape/fixed-tail classification for multipliers
of magnitude at least two. We determine its exceptional basins, retain the
initial division clock, and classify all even multipliers with offset one.
The unit multiplier has an exact ceiling-halving return and logarithmic
capture bound. Mersenne multipliers give a primitive cycle through one with
one affine step and precisely the prescribed number of divisions.
\end{abstract}

\section{Domain and raw clock}

Fix integers \(q,c\), and define
\[
T_{q,c}(n)=
\begin{cases}
n/2,&2\mid n,\\
qn+c,&2\nmid n.
\end{cases}
\]
Division takes precedence, including at zero, so zero is always fixed.
An orbit escapes in magnitude when \(|T^j(n)|\to\infty\) as \(j\to\infty\)
along every raw time, not just a subsequence. All divisibility assertions
in this article are in \(\Z\).

% The odd guard is proved together with the affine formula; no branch is discarded.
\begin{lemma}[Exact odd tail]
\label{lem:tail}
Suppose \(q+c\) is odd. Every nonzero integer has a unique expression
\(n=2^k u\), with \(k\ge0\) and \(u\) odd. The first \(k\) raw steps are
divisions, and for every \(j\ge0\),
\[
T^{k+j}(n)=q^j u+c\sum_{i=0}^{j-1}q^i.
\]
Every state at or after time \(k\) is odd. If \(q\ne1\), then
\begin{equation}
\label{eq:shift}
(q-1)T^{k+j}(n)+c=q^j\bigl((q-1)u+c\bigr).
\end{equation}
\end{lemma}

\begin{proof}
Repeated exact halving terminates on a nonzero integer because absolute
value strictly decreases at each such division. The number of halvings is
\(k=v_2(|n|)\). Uniqueness follows because an equality
\(2^k u=2^\ell v\), with \(k<\ell\), would make \(u\) even.
For odd \(x\), reduction modulo two gives \(qx+c\equiv q+c\equiv1\).
The guard therefore remains odd forever once \(u\) is reached.
Induction in the recurrence \(x_{j+1}=qx_j+c\) gives the displayed sum.
The identity \((q-1)\sum_{i<j}q^i=q^j-1\), proved by telescoping,
gives \eqref{eq:shift}. All these steps have their original raw clock.
\end{proof}

% The vanished shift is the only exceptional case, rather than an assumption of escape.
\begin{theorem}[Signed classification]
\label{thm:signed}
Suppose \(q+c\) is odd and \(|q|\ge2\). For \(n=2^k u\ne0\) as above,
put \(b=(q-1)u+c\).
\begin{enumerate}[label=(\roman*)]
\item If \(b=0\), the orbit is constantly \(u\) after time \(k\).
\item If \(b\ne0\), the full raw orbit escapes in magnitude, and
\[
\lim_{j\to\infty}\frac{|T^{k+j}(n)|}{|q|^j}
=\left|u+\frac{c}{q-1}\right|>0.
\]
\item The only periodic points are zero and the odd integer
\(u_*=-c/(q-1)\), if that quotient is an odd integer.
The basin of \(u_*\) is exactly
\(\{2^k u_*:k\ge0\}\); the basin of zero is \(\{0\}\).
\end{enumerate}
If \(q<0\), every escaping tail eventually alternates signs.
If \(q>0\), its eventual sign is the sign of \(u+c/(q-1)\).
\end{theorem}

\begin{proof}
The integer \(q-1\) is nonzero. If \(b=0\), equation \eqref{eq:shift}
forces every tail state to equal \(u\). Otherwise it gives
\[
T^{k+j}(n)=q^j\left(u+\frac{c}{q-1}\right)-\frac{c}{q-1}.
\]
The first term has magnitude tending to infinity and the second is fixed.
Dividing absolute values by \(|q|^j\) proves the limit. Removing a finite
prefix preserves full-sequence escape. The eventual-sign conclusions
follow from the same formula, with a fixed nonzero leading coefficient.

A periodic nonzero orbit must equal a fixed odd tail: a positive number
of initial divisions cannot belong to an orbit whose later states are
all odd. The equation for that fixed tail is \((q-1)u+c=0\).
Conversely every stated dyadic multiple executes its halving prefix and
then stays fixed. No other odd anchor reaches the fixed tail, since its
nonzero shift is multiplied by nonzero powers of \(q\).
Every nonzero trajectory has only nonzero initial divisions and an odd
tail, so none reaches zero. This proves both basin assertions.
\end{proof}

% Positivity selects the expanding side of the shifted affine line.
\begin{corollary}[Two all-positive escape families]
If \(q\ge2,c\ge1\), \(q\) is even, and \(c\) is odd, every positive
seed escapes to \(+\infty\). The same holds for \(q=3\) and every
nonnegative even offset \(c\).
\end{corollary}

\begin{proof}
The coefficients have opposite parity. For a positive odd anchor \(u\),
\((q-1)u+c>0\) in the first case and \(2u+c>0\) in the second.
Theorem~\ref{thm:signed} applies and gives a positive leading coefficient.
\end{proof}

% The divisibility of one determines the entire even-multiplier exception set.
\begin{corollary}[All even multipliers with offset one]
For every even integer \(q\) with \(|q|\ge2\), the only nonzero
non-escaping seeds are \(-2^k\), \(k\ge0\), when \(q=2\).
They reach the fixed point \(-1\). For all other such even \(q\),
every nonzero signed seed escapes in magnitude. All positive seeds
escape for every even \(q\ge2\).
\end{corollary}

\begin{proof}
An integral exceptional anchor would require \(q-1\mid1\).
The possibilities are \(q=0\) or \(q=2\), and the magnitude assumption
excludes zero. For \(q=2\), the anchor is \(-1\).
The basin and positivity assertions follow from Theorem~\ref{thm:signed}.
\end{proof}

% A target exponent is recovered by decreasing exact quotients, not forward orbit search.
\begin{theorem}[Polynomial-time reachability]
\label{thm:reachability}
Suppose \(q,c\in\Z\), \(q+c\) is odd, and \(|q|\ge2\).
Given signed integers \(n,m\), both reachability of \(m\) from \(n\)
and eventual periodicity of \(n\) are decidable. With binary input,
these decisions take polynomial bit time, uniformly in the coefficients
and both seeds.
\end{theorem}

\begin{proof}
If \(n=0\), its only reachable value is zero, and it is periodic.
Otherwise extract \(n=2^k u\) and check each value \(n/2^i\),
\(0\le i\le k\), for equality with \(m\). This checks the entire
initial division prefix. Put \(b=(q-1)u+c\) and \(e=(q-1)m+c\).
If \(b=0\), the subsequent orbit is fixed at \(u\), which was
already checked. If \(b\ne0\), equation~\eqref{eq:shift} shows that
\(m\) occurs in the tail if and only if
\[
e=q^j b\quad\text{for some integer }j\ge0.
\]
Necessity follows from that equation. Conversely, the equation forces
\(m=T^{k+j}(n)\), since multiplication by the nonzero integer
\(q-1\) is injective. Thus it recovers a legal raw state, not just a
formal affine value.
If \(e=0\), the equation is impossible. Otherwise start at \(r=e\).
Check \(r=b\); if unequal and \(q\mid r\), replace \(r\) by
\(r/q\) and repeat. If \(q\nmid r\), stop with a negative answer.
Each division decreases nonzero magnitude by a factor at least two.
Hence every possible exponent is checked in finitely many steps.
Eventual periodicity is exactly the condition \(n=0\) or \(b=0\),
by Theorem~\ref{thm:signed}.

Let \(L\) be the sum of the binary lengths of \(q,c,n,m\), with
one bit allowed for zero and a separate sign bit. The extraction uses
at most \(L\) halvings. Products forming \(b,e\) have at most
\(O(L)\) bits, and the target-quotient loop uses at most \(O(L)\)
divisions on integers of at most \(O(L)\) bits. Ordinary schoolbook
multiplication and division use \(O(L^2)\) bit operations at these
sizes. Thus \(O(L^3)\) bit operations suffice. No step searches an
unbounded forward orbit. This is a written algorithmic proof, not an
additional Lean declaration.
\end{proof}

\section{The three small multipliers}

% The return is an exact integer operation and has a two-tick ceiling bound.
\begin{theorem}[Unit multiplier capture and its clock]
\label{thm:unit}
For \(q=c=1\), every positive seed reaches one. If \(n\ge1\), it
reaches one in at most \(2\lceil\log_2 n\rceil\) raw steps.
The only positive primitive cycle is \(1,2\).
Every negative seed reaches zero.
\end{theorem}

\begin{proof}
For a positive seed define a return block by
\[
H(n)=\begin{cases}n/2,&n\text{ even},\\(n+1)/2,&n\text{ odd}.
\end{cases}
\]
The even block takes one raw tick; the odd block takes two, because
\(n+1\) is even. In both cases \(H(n)=\lceil n/2\rceil\).
For \(n>1\), this is a smaller positive integer, which already proves
capture by induction. More exactly,
\[
H^h(n)=\left\lceil\frac{n}{2^h}\right\rceil.
\]
To justify induction in this formula, for any real \(x\) and positive
integer \(m\), an integer \(r\) satisfies
\(mr\ge\lceil x\rceil\) if and only if \(mr\ge x\).
Thus \(\lceil\lceil x\rceil/m\rceil=\lceil x/m\rceil\).
Taking \(h=\lceil\log_2 n\rceil\) gives \(H^h(n)=1\), and at most
\(2h\) raw ticks suffice. Since \(1\mapsto2\mapsto1\), any positive
cycle must be this primitive cycle.

If \(n<0\), an even step halves its magnitude, and an odd step sends
\(n\) to \(n+1\). The state stays nonpositive and its absolute value
strictly decreases until zero. Zero is fixed.
\end{proof}

% Constant and reflected affine branches have separate terminal cases.
\begin{theorem}[Constant and reflection multipliers]
For \(q=0,c=1\), every nonzero signed seed reaches the fixed point one.
For \(q=-1,c=1\), every signed seed reaches zero, and every nonzero
seed visits one before reaching zero. Zero is the only periodic point
in the reflection case.
\end{theorem}

\begin{proof}
For \(q=0\), the initial divisions reach an odd integer, which maps to
one; one is odd and maps to itself.
For \(q=-1\), an even nonzero state halves its absolute value. An odd
state \(n\) first maps to the even integer \(1-n\), and then to
\((1-n)/2\). If \(|n|\ge3\), this return has magnitude at most
\((|n|+1)/2<|n|\). Integer descent therefore enters \(\{-1,0,1\}\).
The terminal paths are \(-1\mapsto2\mapsto1\mapsto0\) and
\(1\mapsto0\). An even nonzero state cannot map directly to zero;
the only odd preimage of zero is one. Thus every nonzero path visits
one, and every periodic orbit is the fixed point zero.
\end{proof}

\section{Mersenne excursions}

% Distinct powers of two establish primitivity rather than closure alone.
\begin{theorem}[Primitive cycle through one]
\label{thm:mersenne}
Let \(k\ge1\), \(q=2^k-1\), and \(c=1\). The orbit of one is
\[
1\longmapsto2^k\longmapsto2^{k-1}\longmapsto\cdots
\longmapsto2\longmapsto1.
\]
Its primitive raw period is \(k+1\). Precisely,
\(T^{j+1}(1)=2^{k-j}\) for \(0\le j\le k\).
For \(k=0\), the multiplier is zero and one is fixed.
\end{theorem}

\begin{proof}
The first step is affine and equals \((2^k-1)\cdot1+1=2^k\).
All subsequent states before return are positive even powers of two,
so division is the legal branch. There are exactly \(k\) divisions.
Each state at times \(1,\ldots,k\) is greater than one; hence no earlier
positive time closes the orbit. This proves primitivity. The case
\(k=0\) is the constant multiplier case already computed.
\end{proof}

For example, \(q=3,7,15\) give primitive raw periods \(3,4,5\).
The statement specifies the cycle through one and does not classify
other cycles of odd multipliers.

\section{Formal statements and exact computation}

The companion \texttt{lean/DyadicDynamics.lean} \cite{DyadicCode26}
proves odd guard propagation, the signed division prefix, the integer
shift identity, full-sequence magnitude escape, fixed-tail recovery,
and the signed all-seed dichotomy. Its positive natural-number
\texttt{UnitCapture} theorem proves capture of one for \(q=c=1\),
and its Mersenne theorems prove every phase, closure, and primitive
period. The logarithmic bound and the constant/reflection statements
have the detailed written proofs above; they are not additional Lean
declarations in that companion.
The Go checks use exact integer arithmetic and finite declared ranges;
they test the formulas and clocks, rather than supply the infinite proofs.

\bibliographystyle{amsalpha}
\bibliography{DyadicReferences}
\end{document}
